
\par\noindent\textit{由推论~\ref{cor:7-2-219} 证明命题~\ref{prop:7-2-187}。}
沿用命题~\ref{prop:7-2-187} 的记号。
令 $f_i^{s=0}=f_i^{s=1}=f_i$，
把命题的已知条件纳入定义~\ref{def:7-2-204}
的假设，并核验条件~\ref{cond:7-2-215}；
随后应用推论~\ref{cor:7-2-219}。

取 $\widehat{\mathfrak C}'_i$，
使其两个 $s$-端面
$\widehat{\mathfrak C}'_{i,s=s_0}$
都成为定义~\ref{def:7-2-182bis}
中的 $\mathfrak C_i^{(r)}$。
具体地，依引理 4.2.25，取模型
\[
\widehat{\mathfrak C}'_i
=\bigl(\mathfrak C_i^{(r)}\bigr)^{[0,1]_s}.
\]
相应的 $\operatorname{Eval}_{r=r_0}$、
$\operatorname{Incl}_r$
由 $\mathfrak C_i^{(r)}$ 的映射诱导。
注意，此处 $\mathfrak C_i^{(r)}$
满足 (7.2.189bis)。由于上述取法，
\[
\widehat{\mathfrak C}'_{i,r=0}
=\widehat{\mathfrak C}'_{i,r=1}
=(C'_i)^{[0,1]_s}.
\]

假设已给定
$\mathcal F_i\colon C_i\to
\mathfrak C_{i+1}^{(r)}
=\widehat{\mathfrak C}'_{i+1,s=1}
=\widehat{\mathfrak C}'_{i+1,s=0}$。
令 $\mathcal F_i^{s=0}=\mathcal F_i^{s=1}=\mathcal F_i$；
由 (7.2.186)，这两列映射各自具有
交换误差同伦的性质。
由于 $f_i=f_i^{s=0}=f_i^{s=1}$
都是 $C_i\to C'_i$ 的映射，
可取
\[
\mathcal H_i=\operatorname{Incl}_s\circ f_i
\colon C_i\longrightarrow
(C'_i)^{[0,1]_s}
=\widehat{\mathfrak C}'_{i,r=1}.
\]
其同伦关系显然成立。
按定理 4.2.34，由
$I'_i\colon C'_i\to C'_{i+1}$ 诱导映射
\[
\mathcal I'_i\colon(C'_i)^{[0,1]_s}
\longrightarrow(C'_{i+1})^{[0,1]_s},
\]
它满足 (7.2.206)。最后取
\[
\mathfrak H_i
=\operatorname{Incl}_s\circ\mathcal F_i
\colon C_i\longrightarrow\widehat{\mathfrak C}'_{i+1}.
\]
直接核对即得 (7.2.207)，
从而定义~\ref{def:7-2-204} 的条件全部满足。

再核验条件~\ref{cond:7-2-215}。
$\mathcal I'_i$ 由定理 4.2.34
从 $I'_i$ 诱导，故
\begin{equation}
\mathcal I'_i\circ\operatorname{Incl}_s
=\operatorname{Incl}_s\circ I'_i.
\tag{7.2.220}\label{eq:7-2-220}
\end{equation}
令
\[
\widetilde{\mathfrak C}'_{i+1}
=[0,1]^2\times C'_{i+1}
=\bigl(\mathfrak C'_{i+1}\bigr)^{[0,1]_r},
\qquad
\widetilde{\mathcal H}_i
=\operatorname{Incl}_r\circ
\operatorname{Incl}_s\circ I'_i.
\]
在此情形下 (7.2.189bis) 成立，
因而 (7.2.220) 蕴含 (7.2.216)。
于是，对于 $\{f_i^0\}$ 与 $\{f_i^1\}$
的任意提升 $f^0,f^1$，
推论~\ref{cor:7-2-219} 说明
$f^0$ 与 $f^1$ 滤过 $A_\infty$ 同伦。
命题~\ref{prop:7-2-187} 得证。

\par\noindent\textit{引理~\ref{lem:7-2-208} 的证明。}
设 $C$ 是滤过 $A_\infty$ 代数，
$\{I_i\colon C'_i\to C'_{i+1}\}$
是 $A_{\bar n,\bar k}$ 代数的有向系统。
由定理~\ref{thm:7-2-72} 与
引理~\ref{lem:7-2-69}，
$\{C'_i\}$ 均已提升为滤过 $A_\infty$ 代数。
固定一个 $i$；先证明如下辅助引理。

\begin{lemma}{7.2.221}\label{lem:7-2-221}
设 $\mathfrak g\colon C\to C'_{i+1}$ 是滤过
$A_{\bar n(\lambda_{i+\ell}),\bar k(\lambda_{i+\ell})}$
同态，$\mathfrak f\colon C\to C'_i$ 是滤过
$A_{\bar n(\lambda_{i+\ell-1}),\bar k(\lambda_{i+\ell-1})}$
同态。设 $\mathfrak C_{i+1}^{(r)}$
是 $[0,1]_r\times C'_{i+1}$ 的滤过
$A_{\bar n(\lambda_{i+\ell}),\bar k(\lambda_{i+\ell})}$
代数模型，并设
$\mathcal F\colon C\to\mathfrak C_{i+1}^{(r)}$
是滤过
$A_{\bar n(\lambda_{i+\ell-1}),\bar k(\lambda_{i+\ell-1})}$
同态。假定 $\mathcal F$ 满足下列
滤过同阶同态的恒等式\translatornote{原书写 $I=I_i$，
但 $\mathfrak f:C\to C'_i$ 要求
$I:C'_i\to C'_{i+1}$；证明末尾的
$\mathfrak f$ 取法还与 (7.2.186)
的两个端点评价次序不合。
这些型别与端点冲突保留待核。}：
\begin{align}
\operatorname{Eval}_{r=0}\circ\mathcal F
&=\mathfrak g,
\tag{7.2.222.1}\label{eq:7-2-222-1}\\
\operatorname{Eval}_{r=1}\circ\mathcal F
&=I\circ\mathfrak f,
\qquad I=I_i.
\tag{7.2.222.2}\label{eq:7-2-222-2}
\end{align}
则 $\mathcal F$ 与 $\mathfrak f$ 可延拓为滤过
$A_{\bar n(\lambda_{i+\ell}),\bar k(\lambda_{i+\ell})}$
同态，使 (7.2.222) 在此阶上仍成立。
\end{lemma}

\begin{proof}
可先假定
\[
\bigl((n,K),(n',K')\bigr)
=\Bigl(
\bigl(\bar n(\lambda_{i+\ell-1}),
\bar k(\lambda_{i+\ell-1})\bigr),
\bigl(\bar n(\lambda_{i+\ell}),
\bar k(\lambda_{i+\ell})\bigr)
\Bigr)
\]
为 (7.2.130.1) 或 (7.2.130.2) 的类型之一；
一般情形再由归纳得到。
下面要证明 (7.2.222.1)
在滤过 $A_{n',K'}$ 同态层次上成立。
按定义，对 $\lVert\beta\rVert=n'$，
这写成
\begin{equation}
\operatorname{Eval}_{r=0}\circ
\mathcal F_{K',\beta}
=\mathfrak g_{K',\beta}.
\tag{7.2.223}\label{eq:7-2-223}
\end{equation}
由于 $\mathfrak g$ 是滤过
$A_{n',K'}$ 同态，
$\operatorname{Eval}_{r=0}$ 是
$A_{n',K'}$ 同伦等价，
可先找到满足 (7.2.223) 的映射
$\mathcal F^I_{K',\beta}$；
它未必与已给分量一起构成滤过
$A_{n',K'}$ 同态。
对 $\lVert\beta\rVert=n'$，
令
\[
o_{K',\beta}(\mathcal F)
\in\operatorname{Hom}
\bigl(B_{K'}(\overline C[1]),
\overline{\mathfrak C}_{i+1}^{(r)}[1]\bigr)
\]
为将
$\{\mathcal F_{K,\beta'}\mid
(\lVert\beta'\rVert,K)<(n',K')\}$
延拓为滤过 $A_{n',K'}$ 同态的障碍圈\translatornote{原书将 $\mathfrak K$
定义为两个已移位复形之间映射的核，
但下文又写 $\mathfrak K[1]$；
这里保留源记号，暂不判定是否重复移位。}。
置
\[
\mathfrak K
=\ker\bigl(
\overline{\operatorname{Eval}}_{r=0}
\colon\overline{\mathfrak C}_{i+1}^{(r)}[1]
\longrightarrow\overline C'_{i+1}[1]
\bigr).
\]
于是
\[
o_{K',\beta}(\mathcal F)
+\delta_1(\mathcal F^I_{K',\beta})
\in\operatorname{Hom}
\bigl(B_{K'}(\overline C[1]),\mathfrak K[1]\bigr).
\]
这里
$\delta_1(A)=\overline\delta_1(A)
=\overline{\mathfrak m}_1\circ A
-(-1)^{\deg A}A\circ\widehat{\mathfrak m}_1$。
因 $\overline{\operatorname{Eval}}_{r=0}$
是同伦等价，$\mathfrak K$ 无同调。
故可取
$\operatorname{Cor1}_{K',\beta}\in
\operatorname{Hom}(B_{K'}(\overline C[1]),
\mathfrak K[1])$，使
\[
\mathcal F^{II}_{K',\beta}
:=\mathcal F^I_{K',\beta}
+\operatorname{Cor1}_{K',\beta}
\]
连同所有较低阶分量，定义满足
(7.2.223) 的滤过
$A_{n',K'}$ 同态
$C\to\mathfrak C_{i+1}^{(r)}$。

由引理~\ref{lem:7-2-129}，
可选取 $\mathfrak f^I_{K',\beta}$，
使其与已给低阶分量合成滤过
$A_{n',K'}$ 同态，记作 $\mathfrak f^I$。
由于
\begin{equation}
\overline{\operatorname{Eval}}_{r=1}
\circ\mathcal F^{II}_{K',\beta}
-(I\circ\mathfrak f^I)_{K',\beta}
\in\operatorname{Hom}
\bigl(B_{K'}(\overline C[1]),
\overline C'_{i+1}[1]\bigr)
\tag{7.2.224}\label{eq:7-2-224}
\end{equation}
是 $\delta_1$-圈，
可改变 $\mathfrak f^I_{K',\beta}$，
使 (7.2.224) 成为 $\delta_1$-边界。
亦即，对某个 $\operatorname{Cor2}_{K',\beta}$，
\[
\overline{\operatorname{Eval}}_{r=1}
\circ\mathcal F^{II}_{K',\beta}
-(I\circ\mathfrak f^I)_{K',\beta}
=\delta_1(\operatorname{Cor2}_{K',\beta}).
\]
由 (4.2.2.4) 的满射性，
可取
\[
\operatorname{Cor3}_{K',\beta}
\in\operatorname{Hom}
\bigl(B_{K'}(\overline C[1]),
\overline{\mathfrak C}_{i+1}^{(r)}[1]\bigr)
\]
使
\[
\overline{\operatorname{Eval}}_{r=0}
\circ\operatorname{Cor3}_{K',\beta}=0,
\qquad
\overline{\operatorname{Eval}}_{r=1}
\circ\operatorname{Cor3}_{K',\beta}
=\operatorname{Cor2}_{K',\beta}.
\]
现在令
\[
\mathcal F_{K',\beta}
=\mathcal F^{II}_{K',\beta}
-\delta_1(\operatorname{Cor3}_{K',\beta}).
\]
它具有所要求的性质，引理得证。
\end{proof}

注意，$\widehat{\mathfrak C}'_{i,s=s_0}$
已作为 $[0,1]_r\times C'_i$ 的
滤过 $A_{\bar n(\lambda_i),\bar k(\lambda_i)}$
类模型给定。可将它提升为
$A_\infty$ 类模型；这是命题~\ref{prop:7-2-190}
证明中归纳步的单参数版本。
命题~\ref{prop:7-2-190} 处理双参数情形，
其证明与单参数情形相似而且更不平凡；
故单参数提升过程的细节从略。

还须区分引理~\ref{lem:7-2-208}
与命题~\ref{prop:7-2-217} 的情形。
在命题~\ref{prop:7-2-217} 中，
$\widehat{\mathfrak C}'_{i,r=r_0}$
到滤过 $A_\infty$ 代数的提升已给定，
我们要在该既定结构下，
把 $\operatorname{Incl}_s$ 提升为滤过
$A_\infty$ 同态；见本小节末尾的证明。
而证明引理~\ref{lem:7-2-208} 时，
要同时把
$\widehat{\mathfrak C}'_{i,s=s_0}$
提升为 $A_\infty$ 结构，
把 $\operatorname{Incl}_r$ 提升为
$A_\infty$ 同态。
事实上，引理~\ref{lem:7-2-193} 的证明
就是同时构造 $\mathfrak m_{k,\beta}$
与 $\operatorname{Incl}_{k,\beta}$；
若 $\mathfrak m_{k,\beta}$ 已预先给定，
便不能直接用该证明来构造
$\operatorname{Incl}_{k,\beta}$。

现在完成引理~\ref{lem:7-2-208} 的证明。
对于 $s_0=0,1$，按以下序关系
在 $(\ell,i)$ 上作双重归纳，
把 $f_i^{s=s_0}$ 与
$\mathcal F_i^{s=s_0}$
提升为滤过
$A_{\bar n(\lambda_{i+\ell}),
\bar k(\lambda_{i+\ell})}$ 同态：
\[
(\ell',i')<(\ell,i)
\quad\Longleftrightarrow\quad
\ell'<\ell\ \text{或}\ 
(\ell'=\ell\ \text{且}\ i'<i).
\]
$\ell=0$ 时无需证明。
假定所有小于 $(\ell,i)$ 的情形均已构造，
现考虑 $(\ell,i)$。
记
\[
\bigl(\bar n(\lambda_{i+\ell}),
\bar k(\lambda_{i+\ell})\bigr)
=(n',K').
\]
由于 $(\ell-1,i+1)<(\ell,i)$，
归纳假设分别给出
$\widehat{\mathfrak C}'_{i+1,s=s_0}$
上的滤过 $A_{n',K'}$ 结构，
以及 $A_{n',K'}$ 同态
$f_{i+1}^{s=s_0}$。
还假定 $f_i^{s=s_0}$、
$\mathcal F_i^{s=s_0}$
的滤过 $A_{n,K}$ 同态分量已经给定，
其中 $(n,K)<(n',K')$。

在引理~\ref{lem:7-2-221} 中取
\[
\begin{gathered}
\mathfrak g=f_{i+1}^{s=s_0}\circ I_i,
\qquad
\mathfrak f=I'_i\circ f_i^{s=s_0},
\qquad C=C_i,\\
C'_i=C'_i,\qquad C'_{i+1}=C'_{i+1},
\qquad
\mathfrak C^{(r)}
=\widehat{\mathfrak C}'_{i+1,s=s_0},
\qquad
\mathcal F=\mathcal F_i^{s=s_0}.
\end{gathered}
\]
于是得到
$\widehat{\mathfrak C}'_{i,s=s_0}$
上的滤过 $A_{n',K'}$ 结构，
以及 $f_i^{s=s_0}$、
$\mathcal F_i^{s=s_0}$
的同阶同态分量。
归纳即完成证明。
